\subsection{Crypto Market Microstructure}

\begin{principle} \hlt{Crypto Liquidity Is Exchange-Fragmented}\\
Crypto trades across many centralised and decentralised venues with different custody, margin, liquidation, and fee rules. There is no universal consolidated tape, so fair value and best execution are venue-dependent.
\end{principle}

\begin{taxonomy} \hlt{Crypto Instrument Segmentation}
\begin{enumerate}[label=\roman*.]
\setlength{\itemsep}{0pt}
\item \hlt{spot}: direct token exchange with custody and transfer risk;
\item \hlt{perpetual futures}: leveraged swaps kept near spot through funding;
\item \hlt{dated futures}: basis trades tied to expiry and collateral terms;
\item \hlt{options}: volatility exposure concentrated in fewer venues and strikes;
\item \hlt{on-chain pools}: automated market makers with inventory curves and gas costs.
\end{enumerate}
\end{taxonomy}

\begin{definition} \hlt{Perpetual Funding}\\
Funding is the periodic payment that links perpetual futures to spot. Positive funding means longs pay shorts; negative funding means shorts pay longs. It is both a carry term and a crowding signal.
\end{definition}

\begin{taxonomy} \hlt{Crypto Microstructure Risk Segmentation}
\begin{enumerate}[label=\roman*.]
\setlength{\itemsep}{0pt}
\item \hlt{venue risk}: outage, withdrawal halt, auto-deleveraging, or insolvency;
\item \hlt{liquidation risk}: forced trades from leveraged accounts;
\item \hlt{stablecoin risk}: collateral, redemption, and depeg exposure;
\item \hlt{bridge and settlement risk}: chain congestion, reorgs, and transfer delay;
\item \hlt{fragmentation}: basis and price gaps across venues when capital cannot move quickly.
\end{enumerate}
\end{taxonomy}
